\TieudeT{PHONG TỎA VẬT LÍ 12 - VẬT LÍ NHIỆT}
\subsubsection*{PHẦN I. Thí sinh trả lời câu hỏi từ câu 1 đến câu 18. Mỗi câu hỏi thí sinh chỉ chọn một phương án}
\setcounter{ex}{0}
\begin{ex}
Trong chuyển động nhiệt, các phân tử của chất lỏng
\choice
{chuyển động hoàn toàn hỗn loạn}
{chuyển động hỗn loạn quanh vị trí cân bằng xác định}
{dao động xung quanh các vị trí cân bằng cố định}
{\True dao động quanh vị trí cân bằng nhưng những vị trí cân bằng này không cố định mà di chuyển}
\loigiai{
\begin{itemize}
\item Các phân tử của chất lỏng dao động xung quanh các vị trí cân bằng và các vị trí cân bằng này lại có thể
thay đổi.
\end{itemize}
}
\end{ex}
\begin{ex}
Lực liên kết giữa các phân tử	
\choice
{là lực hút}
{là lực đẩy}
{tuỳ thuộc vào thể của nó, ở thể rắn là lực hút còn ở thể khí là lực đẩy}
{\True gồm cả lực hút và lực đẩy}
\loigiai{
\begin{itemize}
\item Giữa các phân tử có lực hút và lực đẩy, gọi chung là lực liên kết phân tử.
\end{itemize}
}
\end{ex}
\begin{ex}
Thí nghiệm của Brown quan sát
\choice
{hạt phấn hoa bằng kính hiển vi}
{\True chuyển động của hạt phấn hoa trong nước bằng kính hiển vi}
{cánh hoa trong nước bằng kính hiển vi}
{chuyển động của cánh hoa}
\loigiai{
\begin{itemize}
\item Thí nghiệm của Brown quan sát chuyển động của các hạt phấn hoa trong nước, thấy các hạt phấn chuyển động hỗn độn, không ngừng
\end{itemize}
}
\end{ex}
\begin{ex}
Nội năng của một vật
\choice
{chỉ phụ thuộc vào nhiệt độ của vật}
{chỉ phụ thuộc vào thể tích của vật}
{\True phụ thuộc vào thể tích và nhiệt độ của vật}
{không phụ thuộc vào thể tích và nhiệt độ của vật}
\loigiai{
\begin{itemize}
\item Nội năng của một vật phụ thuộc vào nhiệt độ và thể tích của nó.
\item Nhiệt độ liên quan đến động năng của các phân tử, còn thể tích liên quan đến thế năng do tương tác giữa các phân tử.
\end{itemize}
}
\end{ex}
\begin{ex}
Phát biểu nào sau đây nói về nhiệt lượng là \textbf{không} đúng?
\choice
{Nhiệt lượng là số đo độ biến thiên nội năng của vật trong quá trình truyền nhiệt}
{\True Một vật lúc nào cũng có nội năng, do đó lúc nào cũng có nhiệt lượng}
{Đơn vị nhiệt lượng cũng là đơn vị nội năng}
{Nhiệt lượng không phải là nội năng}
\loigiai{
Chỉ khi nào vật có sự biến thiên nội năng do sự truyền nhiệt thì mới có nhiệt lượng.
}
\end{ex}
\begin{ex}
Trong quá trình chất khí tỏa nhiệt và nhận công thì dấu của $A$ và $Q$ trong biểu thức $\Delta U = A+Q$ là
\choice
{$Q > 0,\ A < 0$}
{$Q < 0,\ A < 0$}
{$Q > 0,\ A > 0$}
{\True $Q < 0,\ A > 0$}
\loigiai{
\begin{itemize}
\item Hệ tỏa nhiệt $Q < 0$, nhận công $A > 0$.
\end{itemize}
}
\end{ex}
\begin{ex}
"Độ không tuyệt đối" là nhiệt độ ứng với
\choice
{\True $0 \ K$}
{$0^\circ C$}
{$273^\circ C$}
{$273 \ K$}
\loigiai{
"Độ không tuyệt đối" ứng với nhiệt độ $0 \ K$, nơi mà các phân tử ngừng chuyển động nhiệt hoàn toàn.
}
\end{ex}
\begin{ex}
Nhiệt kế là thiết bị dùng để đo?
\choice
{Chiều dài}
{Thể tích}
{\True Nhiệt độ}
{Diện tích}
\loigiai{
Nhiệt kế là thiết bị dùng để đo nhiệt độ.
}
\haidong{2}\end{ex}
\begin{ex}
Nhiệt kế chất lỏng được chế tạo dựa trên nguyên tắc nào?
\choice
{\True Sự nở vì nhiệt của chất lỏng}
{Sự nở ra của chất lỏng khi nhiệt độ giảm}
{Sự co lại của chất lỏng khi nhiệt độ tăng}
{Sự nở của chất lỏng không phụ thuộc vào nhiệt độ}
\loigiai{
\begin{itemize}
\item Nhiệt kế chất lỏng được chế tạo dựa trên nguyên tắc sự nở vì nhiệt của chất lỏng.
\item Khi nhiệt độ thay đổi, thể tích của chất lỏng trong nhiệt kế sẽ thay đổi tương ứng.
\end{itemize}
}
\haidong{2}\end{ex}
\begin{ex}
Đồ thị nào sau đây biểu diễn mối quan hệ giữa khối lượng nước còn lại trong bình nhiệt lượng kế và thời gian của quá trình hoá hơi của nước?
\choice
{\begin{tikzpicture}[scale=.7,draw=Mapcolor, very thick, >=stealth']
\tkzDefPoints{0/0/o, 3/0/x, 0/3/y}
\tkzDrawSegments[->](o,x o,y)
\tkzLabelPoint[left](o){$O$}
\tkzLabelPoint[above](x){$\tau\ (s)$}
\tkzLabelPoint[above](y){$m\ (kg)$}
\node at (1.5,-0.5) {(1)};
\draw[orange,very thick] (0,1.5) -- (2.75,2.5);
\end{tikzpicture}}
{\begin{tikzpicture}[scale=.7,draw=Mapcolor, very thick, >=stealth']
\tkzDefPoints{0/0/o, 3/0/x, 0/3/y}
\tkzDrawSegments[->](o,x o,y)
\tkzLabelPoint[left](o){$O$}
\tkzLabelPoint[above](x){$\tau\ (s)$}
\tkzLabelPoint[above](y){$m\ (kg)$}
\node at (1.5,-0.5) {(2)};
\draw[orange, very thick] (0,1.5) .. controls (1,1.75) and (2,2) .. (2.75,2.75);
\end{tikzpicture}}
{\True \begin{tikzpicture}[scale=.7,draw=Mapcolor, very thick, >=stealth']
\tkzDefPoints{0/0/o, 3/0/x, 0/3/y}
\tkzDrawSegments[->](o,x o,y)
\tkzLabelPoint[left](o){$O$}
\tkzLabelPoint[above](x){$\tau\ (s)$}
\tkzLabelPoint[above](y){$m\ (kg)$}
\node at (1.5,-0.5) {(3)};
\draw[orange,very thick] (0,1.5) -- (2.5,0.5);
\end{tikzpicture}	}
{\begin{tikzpicture}[scale=.7,draw=Mapcolor, very thick,>=stealth']
\tkzDefPoints{0/0/o, 3/0/x, 0/3/y}
\tkzDrawSegments[->](o,x o,y)
\tkzLabelPoint[left](o){$O$}
\tkzLabelPoint[above](x){$\tau\ (s)$}
\tkzLabelPoint[above](y){$m\ (kg)$}
\node at (1.5,-0.5) {(4)};
\draw[orange,very thick] (0,1.5) -- (2.75,1.5);
\end{tikzpicture}}
\loigiai{Đồ thị 3 vì khối lượng còn lại là hàm bậc nhất của thời gian.
}
\haidong{2}\end{ex}
\begin{ex}
Nhiệt dung riêng của một chất là nhiệt lượng cần truyền cho [1] chất đó để làm cho nhiệt độ của nó tăng thêm $1^\circ C$.
\choice
{\True [1] là $1 \ kg$}
{[1] là $2 \ kg$}
{[1] là $1 \ g$}
{[1] là $2 \ g$}
\loigiai{
Nhiệt dung riêng của một chất là nhiệt lượng cần truyền cho $1 \ kg$ chất đó để làm cho nhiệt độ của nó tăng thêm $1^\circ C$.
}
\end{ex}
\begin{ex}
Đơn vị nhiệt dung riêng là
\choice
{J/kg}
{$kg / J$}
{\True $J/kg.^\circ C$}
{$kg / J.K$}
\loigiai{
Đơn vị của nhiệt dung riêng là J/kg.K, diễn tả nhiệt lượng cần để tăng nhiệt độ của 1 kg chất đó lên 1 K.
}
\end{ex}
\begin{ex}
Biết nhiệt độ ban đầu của nước là $20^\circ C$ và nhiệt dung riêng của nước là $4200 \ J / kg . K$. Để đun sôi 15 kg nước ở áp suất 1 atm cần cung cấp một nhiệt lượng là 
\choice
{\True $5040 \ k J$}
{$5040 \ J$}
{$50,40 \ k J$}
{$5,040 \ J$}
\loigiai{
\begin{itemize}
\item Khối lượng của nước: $m = 15 \ kg$ (vì 1 lít nước tương đương 1 kg)
\item Nước cần tăng nhiệt độ từ $20^\circ C$ lên $100^\circ C$, chênh lệch nhiệt độ $\Delta t = 80^\circ C$
\item Nhiệt lượng cần thiết:
$$Q = mc\Delta t = 15 . 4200 . 80 = 5040000 \ J = 5040 \ kJ$$
\end{itemize}
}
\end{ex}
\begin{ex}
Người ta trộn $1500 \ g$ nước ở $15^\circ C$ với $100 \ g$ nước ở $37^\circ C$. Bỏ qua sự trao đổi nhiệt với môi trường và bình chứa. Nhiệt độ cuối cùng của hỗn hợp là
\choice
{\True $16,375^\circ C$}
{$26,233^\circ C$}
{$52,562^\circ C$}
{$19,852^\circ C$}
\loigiai{
\begin{itemize}
\item Nhiệt lượng $1500 \ g$ nước thu vào:
\begin{align*}
Q_1=m_1 . c .\left(t_2-t_1\right)=1,5 . 4200 .\left(t_2-15\right)
\end{align*}
\item Nhiệt lượng $100 \ g$ nước tỏa ra:
\begin{align*}
Q_2=m_2 . c .\left(t_3-t_2\right)=0,1 . 4200 .\left(37-t_2\right)
\end{align*}
\item Theo phương trình cân bằng nhiệt ta có: $Q_1=Q_2$
\begin{align*}
& \Rightarrow 1,5 . 4200 .\left(t_2-15\right)=0,1 . 4200 .\left(37-t_2\right) \Rightarrow t_2=16,375^\circ C.
\end{align*}
\end{itemize}
}
\end{ex}
\begin{ex}
Đơn vị nào sau đây là đơn vị của nhiệt nóng chảy riêng của vật rắn?
\choice
{Jun trên kilogam độ $(J /kg.K$)}
{\True Jun trên kilogam $(J / kg)$}
{Jun ($J$)}
{Jun trên độ ($J/K$)}
\loigiai{
Đơn vị của nhiệt nóng chảy riêng là Jun trên kilogam $(J/kg)$.
}
\end{ex}
\begin{ex}
Khi vật rắn tinh thể đang nóng chảy thì đại lượng nào của vật \textbf{không} thay đổi?
\choice
{Thể tích của vật}
{Nội năng của vật}
{\True Nhiệt độ của vật}
{Hình dạng của vật}
\loigiai{
\begin{itemize}
\item Trong quá trình nóng chảy của vật rắn tinh thể, nhiệt độ của vật không thay đổi.
\item Điều này có nghĩa là nhiệt độ của vật duy trì ở nhiệt độ nóng chảy suốt quá trình nóng chảy.
\item Lượng nhiệt cung cấp cho vật được sử dụng để phá vỡ cấu trúc tinh thể, chuyển vật từ pha rắn sang pha lỏng mà không làm tăng nhiệt độ.
\item Các đại lượng khác như thể tích và nội năng có thể thay đổi.
\item Thể tích của vật có thể thay đổi do sự nở ra khi chuyển từ rắn sang lỏng, trong khi nội năng tăng do nhiệt lượng hấp thụ.
\end{itemize}
}
\end{ex}
\begin{ex}
Khi nhiệt độ của vật tăng lên thì
\choice
{\True động năng của các phân tử cấu tạo nên vật tăng}
{động năng của các phân tử cấu tạo nên vật giảm}
{nội năng của vật giảm}
{thế năng của các phân tử cấu tạo nên vật tăng}
\loigiai{
\begin{itemize}
\item Khi nhiệt độ của vật tăng lên, năng lượng nhiệt của vật cũng tăng.
\item Điều này đồng nghĩa với việc động năng của các phân tử cấu tạo nên vật tăng lên do các phân tử chuyển động nhanh hơn.
\item Nội năng của vật tăng bởi vì nó bao gồm cả động năng và thế năng của các phân tử.
\end{itemize}
}
\haidong{2}\end{ex}
\begin{ex}
Nhiệt hóa hơi riêng của nước là $2,3 . 10^6\ J/kg$ ở nhiệt độ $100^\circ C$ có nghĩa là
\choice
{một lượng nước bất kì cần thu một lượng nhiệt là $2,3 . 10^6 \ J$ để bay hơi hoàn toàn}
{mỗi kilogam nước cần thu một lượng nhiệt là $2,3 . 10^6 \ J$ để bay hơi hoàn toàn}
{mỗi kilogam nước sẽ tỏa ra một lượng nhiệt là $2,3 . 10^6 \ J$ khi bay hơi hoàn toàn ở nhiệt độ $100^\circ C$}
{\True mỗi kilogam nước cần thu một lượng nhiệt là $2,3 . 10^6 \ J$ để bay hơi hoàn toàn ở nhiệt độ $100^\circ C$}
\loigiai{
\item Nhiệt hóa hơi riêng của nước là giá trị nhiệt lượng cần cung cấp để một kilogam nước bay hơi hoàn toàn ở nhiệt độ sôi ($100^\circ C$) và áp suất tiêu chuẩn ($1\ atm$).
\item  "Một lượng nước bất kì cần thu một lượng nhiệt là $2,3 . 10^6 \ J$ để bay hơi hoàn toàn" không đúng vì nhiệt hóa hơi riêng là nhiệt hóa hơi tính cho 1 đơn vị khối lượng.
\item  "Mỗi kilogam nước cần thu một lượng nhiệt là $2,3 . 10^6 \ J$ để bay hơi hoàn toàn" đúng nhưng thiếu điều kiện cần thiết về nhiệt độ và áp suất.
\item  "Mỗi kilogam nước sẽ tỏa ra một lượng nhiệt là $2,3 . 10^6 \ J$ khi bay hơi hoàn toàn ở nhiệt độ sôi" là sai, vì nước cần thu nhiệt để bay hơi, không tỏa nhiệt.
\item  "Mỗi kilogam nước cần thu một lượng nhiệt là $2,3 . 10^6 \ J$ để bay hơi hoàn toàn ở nhiệt độ sôi và áp suất tiêu chuẩn" là đúng vì nó đầy đủ điều kiện cần thiết.
\item Vậy đáp án đúng là: mỗi kilogam nước cần thu một lượng nhiệt là $2,3 . 10^6 \ J$ để bay hơi hoàn toàn ở nhiệt độ sôi và áp suất tiêu chuẩn.
}
\haidong{3}
\end{ex}
\subsubsection*{PHẦN 2. Thí sinh trả lời từ câu 1 đến câu 4. Trong mỗi ý a), b), c), d) ở mỗi câu, thí sinh chọn đúng hoặc sai.}
\setcounter{ex}{0}
\begin{ex}
\immini[thm]{Hình dưới mô tả mối quan hệ giữa hai thang đo nhiệt độ Kelvin và Fahrenheit ở điều kiện áp suất tiêu chuẩn.}{\begin{tikzpicture}[scale=1,thick,>=stealth']
\tkzDefPoints{0/0/a, 0/4/a', 3/0/b, 3/4/b', 0/2.5/x, 3/2.5/y}
\draw[->] (0,-0.5) -- (0,4.5);
\draw[->] (3,-0.5) -- (3,4.5);
\draw[dashed,thin] (0,2.5) -- (3,2.5);
\draw[dashed,thin] (0,0) -- (3,0);
\draw[dashed,thin] (0,4) -- (3,4);
\node[above left] at (0,4.5) {$T\ (K)$};
\node[above right] at (3,4.5) {$T_F\ (^\circ F)$};
\tkzLabelPoint[left](a){$273$}
\tkzLabelPoint[left](a'){$373$}
\tkzLabelPoint[right](b){$32$}
\tkzLabelPoint[right](b'){$212$}
\node[left] at (0,2.5) {$339$};
\node[above left] at (3,2.5) {$B$};
\node[above right] at (0,2.5) {$A$};
\tkzDrawPoints[size=3,fill=white](a,a',b,b',x,y)
\end{tikzpicture}}
\choiceTF[t]
{Mỗi độ chia trong thang đo nhiệt độ Fahrenheit có độ lớn bằng $1,8$ lần mỗi độ chia trong thang đo nhiệt độ Kelvin}
{\True Độ biến thiên nhiệt độ là $100 \ K$ trên thang đo nhiệt độ Kelvin sẽ tương ứng với độ biến thiên $180^\circ F$ trên thang đo nhiệt độ Fahrenheit}
{\True Tại điểm $B$ trên hình theo thang nhiệt Fahrenheit có nhiệt độ là $150,8{}^\circ F$}
{\True Tại nhiệt độ $574,25$ độ thì giá trị trên hai thang đo là bằng nhau}
\loigiai{
\begin{itemchoice}
\itemch Sai. Mỗi độ chia trong thang đo nhiệt độ Kelvin có độ lớn bằng $1,8$ lần mỗi độ chia trong thang đo nhiệt độ Fahrenheit.
\itemch Đúng. 
\itemch Đúng.
\begin{align*}
\dfrac{339-273}{100}=\dfrac{T_{F}-32}{180} \Rightarrow T_{F}=150,8^\circ F
\end{align*} 
\itemch Đúng. Gọi $x$ là giá trị nhiệt độ bằng nhau trên cả hai thang đo:
\begin{align*}
\dfrac{x-273}{100}=\dfrac{x-32}{180} \Rightarrow x=574,25
\end{align*}
\end{itemchoice}
}
\end{ex}
\begin{ex}
Một lượng khí khi bị nung nóng đã tăng thể tích $0,02\ m^3$ và nội năng biến thiên một lượng $1280\ J$. Biết quá trình trên áp suất không đổi và bằng $2.10^5\ Pa$.
\choiceTFt
{\True Khí bị nung nóng chứng tỏ khối khí nhận được nhiệt lượng}
{\True Quá trình trên khí thực hiện công nên theo qui ước $A<0$}
{Công mà khối khí sinh ra có giá trị là $400\ J$}
{Nhiệt lượng khối khí nhận được là $2720\ J$}
\loigiai{
\begin{itemize}
\item
\end{itemize}
}
\end{ex}
\begin{ex}
Xét quá trình hóa hơi và quá trình sôi của chất lỏng.
\choiceTFt
{\True Sự bay hơi là sự hóa hơi xảy ra ở mặt thoáng của khối chất lỏng}
{Diện tích mặt thoáng càng lớn, tốc độ gió càng lớn thì sự bay hơi diễn ra càng chậm}
{\True Sự hóa hơi xảy ra ở cả mặt thoáng và trong lòng chất của khối chất lỏng khi chất lỏng sôi}
{Áp suất khí trên bề mặt chất lỏng càng lớn thì nhiệt độ sôi của chất lỏng càng nhỏ}
\loigiai{
\begin{itemize}
\item
\end{itemize}
}
\end{ex}
\begin{ex}
Cho đồ thị biểu diễn sự thay đổi nhiệt độ của nước theo thời gian đun như hình bên.
\begin{center} 
\begin{tikzpicture}[draw=Mapcolor, very thick, scale=0.65,transform shape,thick,>=stealth']
\tkzDefPoints{0/0/o, 7/0/x, 0/5.5/y, 0/-1/y', 0/-1/a, 1/0/b, 2/0/c, 5/5/d, 6/5/e}
\draw[help  lines, xstep=0.5cm,ystep=0.5cm]  (0,-1)  grid  (7,5.5);
\tkzDrawSegments[thick,->](o,x y',y)
\tkzLabelPoint[left](o){$O$}
\tkzLabelPoint[above](x){Thời gian (phút)}
\tkzLabelPoint[above](y){Nhiệt độ ($^\circ C$)}
\tkzDrawSegments[very thick, blue](a,b b,c c,d d,e)
\tkzDrawPoints[size=4,fill=white](a,b,c,d,e,o)
\tkzLabelPoint[left](a){$-10$}
\node[left] at (0,5) {$100$};
\foreach  \x  in  {1,2,...,6}
{\pgfmathtruncatemacro{\label}{5*\x};
\node[below](\label) at (\x,0) {\label};}
\end{tikzpicture}
\end{center}
\choiceTFt
{\True Từ phút thứ $0$ đến phút thứ $5$ nước ở thể rắn}
{\True Từ phút thứ $5$ đến phút thứ $10$ xảy ra quá trình nóng chảy}
{Từ phút thứ $10$ đến phút thứ $25$ nước ở thể rắn}
{\True Từ phút thứ $25$ đến phút thứ $30$ xảy ra quá trình sôi}
\end{ex}
\subsubsection*{PHẦN 3. Thí sinh trả lời từ câu 1 đến câu 6.}
\setcounter{ex}{0}
\begin{ex}
Biết nhiệt nóng chảy riêng của thiếc là $0,61.10^5\ J/kg$. Nhiệt lượng cần thiết để làm nóng chảy hoàn toàn khối thiếc là $2,44.10^6\ kJ$. Khối lượng của khối thiếc bằng bao nhiêu tấn?
\SA[oly]{40}
\end{ex}
\begin{ex}
Một người cọ xát một miếng sắt dẹt có khối lượng $200\ g$ trên một tấm đá mài. Sau một thời gian, nhiệt độ của miếng sắt tăng thêm $15^{\circ}C$. Biết nhiệt dung riêng của sắt là $460\ J/kg.K$ và $60\%$ công của người này dùng để làm nóng miếng sắt. Công mà người này đã thực hiện bằng bao nhiêu jun?
\SA[oly]{2300}
\end{ex}
\begin{ex}
Khi truyền nhiệt lượng $3000\ J$ cho một khối khí trong một xilanh hình trụ thì khối khí dãn nở đẩy pit-tông làm thể tích khí tăng thêm $0,005\ m^3$. Giả sử áp suất khối khí trong bình không đổi và bằng $2,4.10^5\ Pa$. Độ biến thiên nội năng của khối khí là bao nhiêu jun?
\SA[oly]{1800}
\end{ex}
\begin{ex}
Một viên bi có khối lượng $m=100 \ g$ được thả từ độ cao $h_1=10 \ m$ xuống mặt sân thì nảy lên đến được độ cao $h_2=4 \ m$, lấy $g=10 \ m/s^2$. Độ biến thiên nội năng của viên bi, mặt sân và không khí là bao nhiêu jun?
\shortans[oly]{6} 
\loigiai{
Độ biến thiên nội năng của hệ viên bi, mặt sân và không khí:
$\Delta U=E_1-E_2=m g\left(h_1-h_2\right)=6\ J$.
}
\end{ex}
\begin{ex}
Các ngôi nhà có tên gọi là Solar Active House là các ngôi nhà được trang bị hệ thống sưởi bằng năng lượng Mặt Trời thông qua bể chứa nước, bộ thu nhiệt và các ống dẫn nước nóng được lắp đặt dưới sàn nhà. Biết nước trong bồn chứa có nhiệt độ $50^\circ C$; nhiệt dung riêng của nước là $c=4180 \ J/kg$.K; khối lượng riêng của nước là $1000 \ kg / m^3$. Nếu trong một ngày mùa đông cần nhiệt lượng $6.10^{8} \ J$ để duy trì nhiệt độ $22^\circ C$ bên trong nhà thì bồn chứa cần có bao nhiêu lít nước (làm tròn kết quả  đến chữ số hàng đơn vị)?
\shortans[oly]{5126} 
\loigiai{
\begin{itemize}
\item Ta có: $t_1=22^\circ C, t_2=50^\circ C$
\item Từ: $Q=mc \Delta t \Rightarrow m=\dfrac{Q}{c \Delta t}=\dfrac{6.10^{8}}{4180.(50-22)}=5126 \ kg$
\item Vậy bồn cần có $5126$ lít nước.
\end{itemize}
}
\end{ex}
\begin{ex}
\immini[thm]{Người ta dùng một lò hồ quang điện để nấu chảy một khối kim loại nặng $30$ kg. Biết lò hồ quang điện mỗi phút cung cấp cho khối kim loại một nhiệt lượng là $400\ kJ$. Đồ thị nhiệt theo thời gian của khối kim loại cho như hình vẽ bên. Nhiệt dung riêng của khối kim loại bằng bao nhiêu $J / k g . K$ (làm tròn kết quả  đến chữ số hàng đơn vị)?
}{
\begin{hinhanh}{7}
\begin{tikzpicture}[draw=Mapcolor, very thick, scale = 1, line join=round, line cap=round, >=stealth,line width=1.5]
\draw[->, line width = 1.5] (0,0)--(0,4) node[right] {$t\ (^0C)$};
\draw[->, line width = 1.5] (0,0)--(5,0) node[right] {$\tau\ (\text{phút})$};
\draw[fill = black] (0,0) circle(2pt) node [below left] {O};
\draw (0,0.5) node [left] {20} -- (1.5,3) -- (3.5,3) -- (4,0);
\draw[dashed,line width=0.8] (0,3) node [left] {1500}-- (1.5,3);
\draw[dashed,line width=0.8] (1.5,3) -- (1.5,0) node [below] {50};
\draw[dashed,line width=0.8] (3.5,3) -- (3.5,0) node [below] {120};
\end{tikzpicture}
\end{hinhanh}}
\shortans[oly]{450}
\end{ex}



